\subsubsection{Measured effect of the pulling subdivision}

The final benchmark runs serially after the suite, audit and publication of
the runtime change. It compares the pinned first constructor and checker at
\code{8354e8aab} against the final pulling implementation, using the same
maintained manifold validator. Each timed call starts afresh, constructs the
exterior, independently replays its source geometry and validates the finite
manifold. Serialization for the stored output hash is outside the timer.
Both variants perform the complete stated operation; they have different
canonical subdivisions and therefore different expected output hashes.

The four inputs are stabilized circles with $1,8,32,128$ crossings. There
are five rounds and four interleaved arms: old, old repeat, new and new
repeat, plus one warmup per arm and input. All eighty measured calls and
sixteen warmups complete in a total driver elapsed time of 24.675 seconds.
Hashes agree within each subdivision across every repeat. The table gives
unpaired time medians and medians of paired ratios.

\begin{center}
\small
\begin{tabular}{rrrrrrr}
\toprule
Crossings & New tets & Old (ms) & New (ms) & Old/new & Old A/A & New A/A \\
\midrule
1 & 20 & 8.444 & 2.269 & 3.748 & 0.996 & 1.043 \\
8 & 160 & 73.140 & 17.259 & 4.346 & 1.032 & 0.997 \\
32 & 640 & 285.732 & 67.210 & 4.228 & 1.021 & 0.998 \\
128 & 2560 & 1217.252 & 284.316 & 4.299 & 0.998 & 0.996 \\
\bottomrule
\end{tabular}

\end{center}

The complete-call old/new ratios are $3.748$--$4.346$. Old A/A ratios are
$0.996$--$1.032$ and new A/A ratios are $0.996$--$1.043$. At 128 crossings,
the median falls from $1.217252$ seconds to $0.284316$ seconds; tetrahedra
fall from 10,240 to 2,560. This is consistent with the exact factor-four
reduction in geometry. It measures construction, source replay and manifold
validation, not a normal-vector search, a simplified minimal exterior or
whole recognition. It supplies no new general quasipolynomial theorem.
